% GENERATED FILE. Edit canonical lessons, then run npm run generate:books.
% canonical-source-hash: fnv1a64:2b0e5b959e28e7e2
% canonical-lessons: ML-W114-dha, ML-W114-ba, ML-W114-chillu-rr, ML-W114-chha, ML-W114-ai-sign, ML-R114-letters-first-pass, ML-R114-letters-second-pass

\chapter{Letters From Days, Family and Body}
\label{ch:ml-letters-days-family-body}
\hlchaptermodality{114}

\begin{chapteropening}
\textbf{By the end of this chapter:} \emph{I can write dha, ba, the chillu r, chha and the ai sign, and I have now written every letter the words of this book use.}

Complete the last lesson of chapter 114: i can write dha, ba, the chillu r, chha and the ai sign, and I have now written every letter the words of this book use.
\end{chapteropening}

\section[dha]{\ml{ധ} (dha) — the letter inside \ml{ബുധൻ}}

\label{lesson:ML-W114-dha}



\begin{quote}
Before the new one: write \ml{ള} and \ml{ഖ} once each, and name them.

Now say \textbf{\ml{ബുധൻ}} (\emph{budhan}), \textquotedblleft{}Wednesday\textquotedblright{}. You have read it for a long time. This lesson takes one letter out of it and puts it in your hand.
\end{quote}

\subsection*{Script you'll notice: \ml{ധ}}
\textbf{\ml{ധ}} — \emph{dha}, a \emph{d} with a puff of breath.

It is inside \textbf{\ml{ബുധൻ}} (\emph{budhan}), \textquotedblleft{}Wednesday\textquotedblright{}, the day named for the planet Mercury.

\subsection*{Writing: \ml{ധ}}
Put your pen on \ml{ധ} and follow its line. Copy the shape you can see — slowly, and larger than it is printed.

\begin{quote}
This book does not yet tell you \textbf{where to start this letter or which way to travel}. It is not written down here until it can be written down with a source. Copying what is in front of you needs no such source, and it is how the shape gets into your hand in the meantime.
\end{quote}

\subsection*{Your turn}
\begin{itemize}
\raggedright
  \item \emph{Look:} at \textbf{\ml{ബുധൻ}} and point at \ml{ധ} inside it
  \item \emph{Trace it:} \ml{ധ} three times, saying \emph{dh} as you finish each one
  \item \emph{Write it:} \ml{ഖ}, then \ml{ധ}, from memory
  \item \emph{Read it:} \textbf{\ml{ബുധൻ}} aloud once more
\end{itemize}

\begin{tcolorbox}[breakable,colback=teal!4,colframe=teal!35!black,title={Before you move on}]
Which letter is this — \ml{ധ}? (\textbf{dha}, \emph{dh}.) Which word did it come from? (\textbf{\ml{ബുധൻ}}, \emph{budhan}, \textquotedblleft{}Wednesday\textquotedblright{}.)
\end{tcolorbox}
\section[ba]{\ml{ബ} (ba) — the letter that opens \ml{ബുധൻ}}

\label{lesson:ML-W114-ba}



\begin{quote}
Before the new one: write \ml{ഖ} and \ml{ധ} once each, and name them.

Now say \textbf{\ml{ബുധൻ}} (\emph{budhan}), \textquotedblleft{}Wednesday\textquotedblright{}. You have read it for a long time. This lesson takes one letter out of it and puts it in your hand.
\end{quote}

\subsection*{Script you'll notice: \ml{ബ}}
\textbf{\ml{ബ}} — \emph{ba}, the sound \emph{b}.

It opens \textbf{\ml{ബുധൻ}} (\emph{budhan}), the word \textbf{\ml{ധ}} came from in the last lesson.

\subsection*{Writing: \ml{ബ}}
Put your pen on \ml{ബ} and follow its line. Copy the shape you can see — slowly, and larger than it is printed.

\begin{quote}
This book does not yet tell you \textbf{where to start this letter or which way to travel}. It is not written down here until it can be written down with a source. Copying what is in front of you needs no such source, and it is how the shape gets into your hand in the meantime.
\end{quote}

\subsection*{Your turn}
\begin{itemize}
\raggedright
  \item \emph{Look:} at \textbf{\ml{ബുധൻ}} and point at \ml{ബ} inside it
  \item \emph{Trace it:} \ml{ബ} three times, saying \emph{b} as you finish each one
  \item \emph{Write it:} \ml{ധ}, then \ml{ബ}, from memory
  \item \emph{Read it:} \textbf{\ml{ബുധൻ}} aloud once more
\end{itemize}

\begin{tcolorbox}[breakable,colback=teal!4,colframe=teal!35!black,title={Before you move on}]
Which letter is this — \ml{ബ}? (\textbf{ba}, \emph{b}.) Which word did it come from? (\textbf{\ml{ബുധൻ}}, \emph{budhan}, \textquotedblleft{}Wednesday\textquotedblright{}.)
\end{tcolorbox}
\section[chillu r]{\ml{ർ} (chillu r) — the r that ends \ml{ഞായർ}}

\label{lesson:ML-W114-chillu-rr}



\begin{quote}
Before the new one: write \ml{ധ} and \ml{ബ} once each, and name them.

Now say \textbf{\ml{ഞായർ}} (\emph{ñāyar}), \textquotedblleft{}Sunday\textquotedblright{}. You have read it for a long time. This lesson takes one letter out of it and puts it in your hand.
\end{quote}

\subsection*{Script you'll notice: \ml{ർ}}
\textbf{\ml{ർ}} — \emph{chillu r}, an \emph{r} with no vowel after it.

It ends \textbf{\ml{ഞായർ}} (\emph{ñāyar}), \textquotedblleft{}Sunday\textquotedblright{}. A \emph{chillu} is the form a consonant takes when it closes a word with no vowel of its own. You already write other chillus; this is the \emph{r} one.

\subsection*{Writing: \ml{ർ}}
\hlblockfigure{\detokenize{figures/ML-W114-chillu-rr-filmstrip.pdf}}{How \ml{ർ} is written, stroke by stroke}

\begin{itemize}
\raggedright
  \item \textbf{1.} start here: climb around the left arch and carry the upper shoulder right
  \item \textbf{2.} without lifting, sweep clockwise around the right loop and return to the upper crossing
  \item \textbf{3.} without lifting, rise into the chillu hook and curl left above the line — and only now lift
\end{itemize}

\textbf{Pen lifts: 0.} The pen never leaves the paper.

\begin{quote}
Stroke order is one attested teaching order; hands and teachers vary. Source: Sriveenkat, Ml \ml{ർ} order.gif, Malayalam chillu RR stroke-order animation, 00:03.0–00:08.5 (Wikimedia Commons, 2 July 2023).
\end{quote}

\subsection*{Your turn}
\begin{itemize}
\raggedright
  \item \emph{Look:} at \textbf{\ml{ഞായർ}} and point at \ml{ർ} inside it
  \item \emph{Trace it:} \ml{ർ} three times, saying \emph{r} as you finish each one
  \item \emph{Write it:} \ml{ബ}, then \ml{ർ}, from memory
  \item \emph{Read it:} \textbf{\ml{ഞായർ}} aloud once more
\end{itemize}

\begin{tcolorbox}[breakable,colback=teal!4,colframe=teal!35!black,title={Before you move on}]
Which letter is this — \ml{ർ}? (\textbf{chillu r}, \emph{r}.) Which word did it come from? (\textbf{\ml{ഞായർ}}, \emph{ñāyar}, \textquotedblleft{}Sunday\textquotedblright{}.)
\end{tcolorbox}
\section[chha]{\ml{ഛ} (chha) — the letter inside \ml{അച്ഛൻ}}

\label{lesson:ML-W114-chha}



\begin{quote}
Before the new one: write \ml{ബ} and \ml{ർ} once each, and name them.

Now say \textbf{\ml{അച്ഛൻ}} (\emph{acchan}), \textquotedblleft{}father\textquotedblright{}. You have read it for a long time. This lesson takes one letter out of it and puts it in your hand.
\end{quote}

\subsection*{Script you'll notice: \ml{ഛ}}
\textbf{\ml{ഛ}} — \emph{chha}, a \emph{ch} with a puff of breath.

It is inside \textbf{\ml{അച്ഛൻ}} (\emph{acchan}), \textquotedblleft{}father\textquotedblright{}, joined under \textbf{\ml{ച}} in \textbf{\ml{ച്ഛ}}. You have been reading it inside that join since the family chapter.

\subsection*{Writing: \ml{ഛ}}
Put your pen on \ml{ഛ} and follow its line. Copy the shape you can see — slowly, and larger than it is printed.

\begin{quote}
This book does not yet tell you \textbf{where to start this letter or which way to travel}. It is not written down here until it can be written down with a source. Copying what is in front of you needs no such source, and it is how the shape gets into your hand in the meantime.
\end{quote}

\subsection*{Your turn}
\begin{itemize}
\raggedright
  \item \emph{Look:} at \textbf{\ml{അച്ഛൻ}} and point at \ml{ഛ} inside it
  \item \emph{Trace it:} \ml{ഛ} three times, saying \emph{chh} as you finish each one
  \item \emph{Write it:} \ml{ർ}, then \ml{ഛ}, from memory
  \item \emph{Read it:} \textbf{\ml{അച്ഛൻ}} aloud once more
\end{itemize}

\begin{tcolorbox}[breakable,colback=teal!4,colframe=teal!35!black,title={Before you move on}]
Which letter is this — \ml{ഛ}? (\textbf{chha}, \emph{chh}.) Which word did it come from? (\textbf{\ml{അച്ഛൻ}}, \emph{acchan}, \textquotedblleft{}father\textquotedblright{}.)
\end{tcolorbox}
\section[ai sign]{\ml{◌ൈ} (ai sign) — the ai sign in \ml{കൈ}}

\label{lesson:ML-W114-ai-sign}



\begin{quote}
Before the new one: write \ml{ർ} and \ml{ഛ} once each, and name them.

Now say \textbf{\ml{കൈ}} (\emph{kai}), \textquotedblleft{}hand\textquotedblright{}. You have read it for a long time. This lesson takes one letter out of it and puts it in your hand.
\end{quote}

\subsection*{Script you'll notice: \ml{◌ൈ}}
\textbf{\ml{ൈ}} — the \emph{ai} sign.

It is in \textbf{\ml{കൈ}} (\emph{kai}), \textquotedblleft{}hand\textquotedblright{}. Like \textbf{\ml{െ}} and \textbf{\ml{േ}}, which you already read, it is written \textbf{before} the consonant and said \textbf{after} it.

\subsection*{Writing: \ml{◌ൈ}}
Put your pen on \ml{◌ൈ} and follow its line. Copy the shape you can see — slowly, and larger than it is printed.

\begin{quote}
This book does not yet tell you \textbf{where to start this letter or which way to travel}. It is not written down here until it can be written down with a source. Copying what is in front of you needs no such source, and it is how the shape gets into your hand in the meantime.
\end{quote}

\subsection*{Your turn}
\begin{itemize}
\raggedright
  \item \emph{Look:} at \textbf{\ml{കൈ}} and point at \ml{◌ൈ} inside it
  \item \emph{Trace it:} \ml{◌ൈ} three times, saying \emph{ai} as you finish each one
  \item \emph{Write it:} \ml{ഛ}, then \ml{◌ൈ}, from memory
  \item \emph{Read it:} \textbf{\ml{കൈ}} aloud once more
\end{itemize}

\begin{tcolorbox}[breakable,colback=teal!4,colframe=teal!35!black,title={Before you move on}]
Which letter is this — \ml{◌ൈ}? (\textbf{ai sign}, \emph{ai}.) Which word did it come from? (\textbf{\ml{കൈ}}, \emph{kai}, \textquotedblleft{}hand\textquotedblright{}.)
\end{tcolorbox}
\section[onpatu akṣarangaḷ]{(nine letters) — From the letter to the word — a first pass}

\label{lesson:ML-R114-letters-first-pass}



\begin{quote}
Say \textbf{\ml{ശരി}} (\emph{śari}), then write the letter it gave you.
\end{quote}

\subsection*{Your turn}
\noindent
\begin{tabularx}{\linewidth}{@{}>{\raggedright\arraybackslash}X>{\raggedright\arraybackslash}X>{\raggedright\arraybackslash}X@{}}
\toprule
\textbf{letter} & \textbf{name} & \textbf{the word it came from} \\
\midrule
\ml{ശ} & \emph{śa} & \textbf{\ml{ശരി}} \emph{śari} \\
\ml{ങ} & \emph{ṅa} & \textbf{\ml{നിങ്ങൾ}} \emph{niṅṅaḷ} \\
\ml{ള} & \emph{ḷa} & \textbf{\ml{നാളെ}} \emph{nāḷe} \\
\ml{ഖ} & \emph{kha} & \textbf{\ml{സുഖം}} \emph{sukham} \\
\ml{ധ} & \emph{dha} & \textbf{\ml{ബുധൻ}} \emph{budhan} \\
\ml{ബ} & \emph{ba} & \textbf{\ml{ബുധൻ}} \emph{budhan} \\
\ml{ർ} & \emph{chillu r} & \textbf{\ml{ഞായർ}} \emph{ñāyar} \\
\ml{ഛ} & \emph{chha} & \textbf{\ml{അച്ഛൻ}} \emph{acchan} \\
\ml{◌ൈ} & \emph{ai sign} & \textbf{\ml{കൈ}} \emph{kai} \\
\bottomrule
\end{tabularx}

\begin{itemize}
\raggedright
  \item \emph{Write it:} each letter in the table from memory, then check it
  \item \emph{Say it:} the word each one came from
  \item \emph{Read it:} each word in the table aloud, and point at its letter
\end{itemize}

\begin{tcolorbox}[breakable,colback=teal!4,colframe=teal!35!black,title={Before you move on}]
Which word did each letter come from? (\textbf{\ml{ശ}} from \textbf{\ml{ശരി}}; \textbf{\ml{ങ}} from \textbf{\ml{നിങ്ങൾ}}; \textbf{\ml{ള}} from \textbf{\ml{നാളെ}}; \textbf{\ml{ഖ}} from \textbf{\ml{സുഖം}}; \textbf{\ml{ധ}} from \textbf{\ml{ബുധൻ}}; \textbf{\ml{ബ}} from \textbf{\ml{ബുധൻ}}; \textbf{\ml{ർ}} from \textbf{\ml{ഞായർ}}; \textbf{\ml{ഛ}} from \textbf{\ml{അച്ഛൻ}}; \textbf{\ml{◌ൈ}} from \textbf{\ml{കൈ}}.)
\end{tcolorbox}
\section[onpatu akṣarangaḷ]{(nine letters) — From the word back to the letter — a second pass}

\label{lesson:ML-R114-letters-second-pass}



\begin{quote}
Write \textbf{\ml{◌ൈ}} without looking, then say the word it came from.
\end{quote}

\subsection*{Your turn}
\noindent
\begin{tabularx}{\linewidth}{@{}>{\raggedright\arraybackslash}X>{\raggedright\arraybackslash}X>{\raggedright\arraybackslash}X@{}}
\toprule
\textbf{letter} & \textbf{name} & \textbf{the word it came from} \\
\midrule
\ml{ശ} & \emph{śa} & \textbf{\ml{ശരി}} \emph{śari} \\
\ml{ങ} & \emph{ṅa} & \textbf{\ml{നിങ്ങൾ}} \emph{niṅṅaḷ} \\
\ml{ള} & \emph{ḷa} & \textbf{\ml{നാളെ}} \emph{nāḷe} \\
\ml{ഖ} & \emph{kha} & \textbf{\ml{സുഖം}} \emph{sukham} \\
\ml{ധ} & \emph{dha} & \textbf{\ml{ബുധൻ}} \emph{budhan} \\
\ml{ബ} & \emph{ba} & \textbf{\ml{ബുധൻ}} \emph{budhan} \\
\ml{ർ} & \emph{chillu r} & \textbf{\ml{ഞായർ}} \emph{ñāyar} \\
\ml{ഛ} & \emph{chha} & \textbf{\ml{അച്ഛൻ}} \emph{acchan} \\
\ml{◌ൈ} & \emph{ai sign} & \textbf{\ml{കൈ}} \emph{kai} \\
\bottomrule
\end{tabularx}

\begin{itemize}
\raggedright
  \item \emph{Say it:} each word in the table, then write the letter it gave you without looking
  \item \emph{Write it:} each letter once more, slowly, larger than it is printed
\end{itemize}

\begin{tcolorbox}[breakable,colback=teal!4,colframe=teal!35!black,title={Before you move on}]
Say each word in the table once more, and name the letter you took from it.
\end{tcolorbox}
